%%%%%%%%%%%%%%%%%%%%%%%%%%%%%%%%%%%%%%%%%%%%%%%%%%%%%%%%%%%%%%%%%%%%%%
% Lớp 10 | BPT và hệ BPT bậc nhất hai ẩn | Bất phương trình bậc nhất hai ẩn
% Dạng 0D2?1-1: Các khái niệm về BPT bậc nhất hai ẩn
% Số câu: 50  (sinh tự động bởi tach_cau.py; từ nay thầy sửa trực tiếp file này)
%%%%%%%%%%%%%%%%%%%%%%%%%%%%%%%%%%%%%%%%%%%%%%%%%%%%%%%%%%%%%%%%%%%%%%

%%%%%%%%%%%%%%%% TRẮC NGHIỆM 4 PHƯƠNG ÁN %%%%%%%%%%%%%%%%

% nguồn: MaHoaNH_TracNghiem_Lop10
\begin{ex}%[0D2N1-1]%[ma:07a60959]
	Cặp số nào sau đây \textbf{không} là nghiệm của  bất phương trình $x+2y-3>0$.
	 \choice
		 {$\left(-2;3\right)$}
		{$\left(-1;4\right)$} 
		 {$\left(4;0\right)$}
		{\True$\left(-1;0\right)$} 
	 \loigiai{Ta có $-1+2 \cdot 0-3 < 0$ nên cặp số $\left(-1;0\right)$ không là nghiệm của bất phương trình $x+2y-3>0$.}
\end{ex}

% nguồn: MaHoaNH_TracNghiem_Lop10
\begin{ex}%[0D2N1-1]%[ma:ab7e55f8]
	Tìm cặp số  là nghiệm của bất phương trình $-2x+3y-5>0$.
	\choice 
		{$\left(-1;1\right)$} 
		{\True $\left(-1;2\right)$} 
		{$\left(0;0\right)$} 
		{$\left(3;1\right)$} 
	\loigiai{Ta có $-2 \cdot (-1)+3 \cdot 2-5>0$ nên cặp số $\left(-1;2\right)$  là nghiệm của bất phương trình $-2x+3y-5>0$.}
\end{ex}

% nguồn: MaHoaNH_TracNghiem_Lop10
\begin{ex}%[0D2N1-1]%[ma:e6805bd5]
	Cặp số nào sau đây \textbf{không} là nghiệm của bất phương trình  $2x+y-7>0$.
	 \choice
		 {$\left(3;2\right)$}
		 {$\left(5;-1\right)$}
		 {$\left(4;0\right)$}
		{\True $\left(-2;5\right)$} 
	\loigiai{Ta có $2 \cdot (-2)+5-7<0$ nên  $\left(-2;5\right)$ không là nghiệm của bất phương trình $2x+y-7>0$.}
\end{ex}

% nguồn: MaHoaNH_TracNghiem_Lop10
\begin{ex}%[0D2N1-1]%[ma:30abb518]
	Điểm $A(-1;3)$ là điểm \textbf{không} thuộc  miền nghiệm của bất phương trình nào sau đây?
	\choice 
		 {\True $-3x+2y-4<0$}
		{$x+3y>0$} 
		{$3x-y \le 0$} 
		 {$2x-y+4<0$}
	 \loigiai{Vì $-3 \cdot (-1)+2 \cdot 3-4>0$ nên điểm $A(-1;3)$ là điểm không thuộc miền nghiệm của bất phương trình $-3x+2y-4<0$.}
\end{ex}

% nguồn: MaHoaNH_TracNghiem_Lop10
\begin{ex}%[0D2N1-1]%[ma:9ab08552]
	Cặp số nào  sau đây là nghiệm của bất phương trình $2x-y+1<0$?
	\choice 
		{$\left(0;-1\right)$} 
		{$\left(3;5\right)$} 
		{\True $\left(1;4\right)$} 
		{$\left(2;-1\right)$} 
	\loigiai{Ta có $2 \cdot 1-4+1<0$ nên $\left(1;4\right)$ là  nghiệm của bất phương trình $2x-y+1<0$.}
\end{ex}

% nguồn: MaHoaNH_TracNghiem_Lop10
\begin{ex}%[0D2N1-1]%[ma:7b09180c]
	Cặp số $(x;y)=(19;8)$ là một nghiệm của bất  phương trình nào sau đây?
	 \choice
		 {$2x-5y-1>0$}
		 {$x-y<0$}
		 {$3x-y-5<0$}
		{\True  $x-2y+1 \ge 0$}
	\loigiai{Ta có $19 - 2 \cdot 8+1 \ge 0$ nên cặp số $(x;y)=(19;8)$ là một nghiệm  của bất phương trình $x-2y+1 \ge 0$.}
\end{ex}

% nguồn: MaHoaNH_TracNghiem_Lop10
\begin{ex}%[0D2N1-1]%[ma:aa4e9dcb]
	Bất phương trình nào sau đây là bất phương trình bậc nhất hai  ẩn?
	\choice 
		 {$2x^2+3y>0$}
		{$x^2+y^2<2$} 
		{$x+y^2\ge 0$} 
		{\True $2x+3y\ge 0$} 
	 \loigiai{
	Bất phương trình bậc nhất hai ẩn $x$, $y$ là bất phương trình  có một trong các dạng $ax+by>c$; $ax+by\ge c$; $ax+by<c$; $ax+by\le c$ $(a^2+b^2>0)$.}
\end{ex}

% nguồn: MaHoaNH_TracNghiem_Lop10
\begin{ex}%[0D2N1-1]%[ma:a83a5d82]
	Trong các bất phương trình sau, bất phương trình nào là bất phương trình bậc nhất hai ẩn? 
	 \choice
		{$3x^2-x+4\le  0$}
		 {\True $2x+y < 5$}
		{$x^2+5y > 1$} 
		{$2x-5y+z\ge  0$}
	 \loigiai{
	Theo định nghĩa bất phương trình bậc nhất hai ẩn suy  ra.}
\end{ex}

% nguồn: MaHoaNH_TracNghiem_Lop10
\begin{ex}%[0D2N1-1]%[ma:6f617d49]
	Bất phương trình nào sau đây là bất phương trình bậc nhất hai  ẩn?
	\choice 
		 {\True $7x+3y\ge 6$}
		 {$5y^3-3<7$}
		 {$(7x-6y)(3x+7y)\le 7$}
		{$6x^2+2y^2<2$} 
	\loigiai{ 
		Bất phương trình bậc nhất hai ẩn là  $7x+3y\ge 6$.
	} 
\end{ex}

% nguồn: MaHoaNH_TracNghiem_Lop10
\begin{ex}%[0D2N1-1]%[ma:e7932971]
	Bất phương trình  nào sau đây là bất phương trình bậc nhất hai ẩn?
	\choice 
		 {$x+y-2xy > 0$}
		{$x^2-2x+1<  0$}
		 {\True $x-2y-3> 0$}
		{$ax+by-c >  0$}
	\loigiai{ 
		$x-2y-3> 0$ là bất  phương trình bậc nhất hai ẩn.
	} 
\end{ex}

% nguồn: MaHoaNH_TracNghiem_Lop10
\begin{ex}%[0D2N1-1]%[ma:559c354c]
	Cặp  số $\left(2;-1\right)$ là nghiệm của bất phương trình nào dưới đây?
	\choice 
		{$x+y-1< 0$} 
		{$2x+3y >  1$}
		 {$x+3y\ge 0$}
		{\True $x+2y\le  0$}
	\loigiai{ 
		Lần lượt  thay $x=2$; $y=-1$ vào bất phương trình $x+2y\le 0$ ta thấy thỏa mãn.
	} 
\end{ex}

% nguồn: MaHoaNH_TracNghiem_Lop10
\begin{ex}%[0D2N1-1]%[ma:29cfd064]
	Trong các bất  phương trình sau đây, bất phương trình nào là bất phương trình bậc nhất hai ẩn?
	\choice 
		{\True  $2x-3y\ge 1$}
		{$x-2xy\ge  1$}
		{$2x^2-y >  0$}
		{$x^2-2y^3>  3$}
	 \loigiai{
		$2x-3y\ge 1$ là bất  phương trình bậc nhất hai ẩn.
	} 
\end{ex}

% nguồn: MaHoaNH_TracNghiem_Lop10
\begin{ex}%[0D2N1-1]%[ma:0706b424]
	Cặp số $(1;-1)$ là nghiệm của bất phương trình nào sau đây? 
	\choice 
		 {$x+y-3> 0$}
		{$-x-y <  0$}
		{\True  $x+3y+1< 0$}
		{$-x-3y-1<  0$}
	 \loigiai{
		$(1;-1)$ là nghiệm của bất phương  trình $x+3y+1< 0$.
	} 
\end{ex}

% nguồn: MaHoaNH_TracNghiem_Lop10
\begin{ex}%[0D2N1-1]%[ma:9ba17bf4]
	Trong các bất phương trình sau đây, bất phương trình nào là bất phương trình bậc nhất  hai ẩn?
	\choice 
		{$2x^2-y >  0$}
		 {$x^2-2y^3 > 3$}
		{\True  $2x-y < 3$}
		{$x-2y^2\le 1$} 
	 \loigiai{
		$2x-y < 3$ là bất phương trình bậc  nhất hai ẩn.
	 }
\end{ex}

% nguồn: MaHoaNH_TracNghiem_Lop10
\begin{ex}%[0D2N1-1]%[ma:01e50546]
	Bất phương trình nào sau đây là bất phương trình bậc nhất  hai ẩn?
	 \choice
		{$x^2+y-2\le  0$}
		 {$y^2> 0$}
		{$x\left(y-1\right)\le  0$}
		{\True $x+3y+2> 0$} 
	 \loigiai{
		$x+3y+2> 0$ là  bất phương trình bậc nhất hai ẩn.
	 }
\end{ex}

% nguồn: MaHoaNH_TracNghiem_Lop10
\begin{ex}%[0D2N1-1]%[ma:beeeb709]
	Trong các bất phương trình sau,  bất phương trình nào \textbf{không} phải là bất phương trình bậc nhất hai ẩn?
	 \choice
		 {$2x-3y-2022\le 0$}
		{$5x+y-2022\ge 2x+11$} 
		{$x+2025>0$} 
		{\True $\dfrac{x}{y}+1>0$} 
	 \loigiai{
		Bất phương trình bậc nhất hai ẩn có một trong các dạng $ax + y <c$; $ax + y >c$; $ax + y \le c$; $ax + y \ge c$, trong đó $a$,  $b$, $c$ là những số cho trước, $a$ và $b$ không đồng thời bằng $0$, $x$ và $y$ là các ẩn.\\
		Do đó $\dfrac{x}{y}+1>0$ không phải là bất phương trình bậc nhất hai ẩn. 
	} 
\end{ex}

% nguồn: MaHoaNH_TracNghiem_Lop10
\begin{ex}%[0D2N1-1]%[ma:806f2684]
	Trong  các bất phương trình sau, bất phương trình nào là bất phương trình bậc nhất hai ẩn?
	\choice 
		{$2x^2+3y>4$} 
		 {$xy+x<5$}
		{\True $3^2x+4^3y\ge 6$} 
		{$x+y^3\le 3$} 
	\loigiai{ 
		Bất phương trình bậc nhất hai ẩn có một trong các dạng $ax + y <c$; $ax + y >c$; $ax + y \le c$; $ax  + y \ge c$, trong đó $a$, $b$, $c$ là những số cho trước, $a$ và $b$ không đồng thời bằng $0$, $x$ và $y$ là các ẩn.\\
		Do đó $3^2x+4^3y\ge  6\Leftrightarrow 9x+64y\ge 6$.
	 }
\end{ex}

% nguồn: MaHoaNH_TracNghiem_Lop10
\begin{ex}%[0D2N1-1]%[ma:d7e5708e]
	Với giá trị nào của tham số $m$ thì  bất phương trình $\left(m^2-1\right)x+\left(2-2m\right)y>2$ là bất phương trình bậc nhất hai ẩn.
	\choice 
		 {\True $m \neq 1$}
		 {$m \neq -1$}
		 {$m=1$}
		 {$m=-1$}
	\loigiai{Bất phương trình $\left(m^2-1\right)x+\left(2-2m\right)y>2$ là bất phương trình  bậc nhất hai ẩn khi và chỉ khi
		$$(m^2-1)^2+(2-2m)^2 \neq 0 \Leftrightarrow (m-1)^2(m+1)^2+4(m-1)^2 \neq 0  \Leftrightarrow (m-1)^2 \left((m+1)^2+4\right) \neq 0.$$
		Ta có $(m+1)^2+4 \ge 4>0$, $\forall m \in  \mathbb{R}$ nên
	$$(m-1)^2 \left((m+1)^2+4\right) \neq  0 \Leftrightarrow (m-1)^2 \neq 0 \Leftrightarrow m=1.$$}
\end{ex}

% nguồn: MaHoaNH_TracNghiem_Lop10
\begin{ex}%[0D2N1-1]%[ma:b2b2fda1]
	Miền nghiệm của bất phương trình $3x+2\left(y+3\right) >  4\left(x+1\right)-y+3$ là nửa mặt phẳng chứa điểm:
	\choice 
		 {$\left(3; 0\right)$}
		{$\left(3;  1\right)$}
		 {\True $\left(1; 2\right)$}
		{$\left(0;  1\right)$}
	\loigiai{ 
		Ta có $3x+2\left(y+3\right) >  4\left(x+1\right)-y+3\Leftrightarrow-x+3y-1> 0$.\\
		Vì $-1+3\cdot 2-1> 0$ là mệnh đề đúng nên miền nghiệm  của bất phương trình trên chứa điểm có tọa độ $\left(1; 2\right)$.
		} 
\end{ex}

% nguồn: MaHoaNH_TracNghiem_Lop10
\begin{ex}%[0D2H1-1]%[ma:0e83c15c]
	Tìm  $m$ để bất phương trình $mx-3y<2$ là bất phương trình bậc nhất hai ẩn?
	 \choice
		{\True  $m \in \mathbb{R}$}
		{$m  \neq 0$}
		{$m>0$} 
		{$m<0$} 
	\loigiai{Bất phương  trình $mx-3y<2$ là bất phương trình bậc nhất hai ẩn khi và chỉ khi
		$$m^2+3 \neq 0 \Leftrightarrow  m^2 \neq -3 \; \; (\text{luôn đúng}).$$
	 Vậy bất phương trình $mx-3y<2$ là bất phương trình bậc nhất hai ẩn với mọi $m \in \mathbb{R}$.}
\end{ex}

% nguồn: MaHoaNH_TracNghiem_Lop10
\begin{ex}%[0D2H1-1]%[ma:0bdd4d2a]
	Bất phương trình nào sau đây  là bất phương trình bậc nhất hai ẩn?
	\choice 
		{$\dfrac{-1}{7}x-\dfrac{y^2}{3} \le  8$ }
		{$\sqrt{2} x^2-5 \sqrt{y}  \ge 8$}
		{$2  \cdot \dfrac{1}{x}-5\cdot \dfrac{1}{y}>8$}
		{\True $-\dfrac{2}{5} x-5^2  y \leq-\sqrt{15}$}
	 \loigiai{
		Bất phương trình bậc nhất hai ẩn có một trong các dạng $ax + y <c$; $ax + y >c$; $ax + y \le c$;  $ax + y \ge c$, trong đó $a$, $b$, $c$ là những số cho trước, $a$ và $b$ không đồng thời bằng $0$, $x$ và $y$ là các ẩn.\\
		Do đó bất phương trình $-\dfrac{2}{5} x-5^2 y \leq-\sqrt{15}$ là bất phương trình  bậc nhất hai ẩn $x$, $y$.
	} 
\end{ex}

% nguồn: MaHoaNH_TracNghiem_Lop10
\begin{ex}%[0D2H1-1]%[ma:b49ff26e]
	Với giá trị nào của tham số $m$  thì bất phương trình $(m^2 - 3m + 2)x + (1 - m)y \ge 3$ là bất phương trình bậc nhất hai ẩn?
	\choice 
		{\True $m \ne 1$} 
		{$\heva{& m \ne  1 \\ &m \ne 2 }$}
		 {$m = 1$}
		 {$m = 2$}
	\loigiai{ 
		Để $(m^2 - 3m +  2)x + (1 - m)y \ge 3$ là bất phương trình bậc nhất hai ẩn thì
		\[(m^2 - 3m + 2)^2 + (1 - m)^2 \ne 0 \Leftrightarrow (m - 1)^2\left[(m - 2)^2 + 1\right] \ne 0 \Leftrightarrow m  \ne 1.\]
	 }
\end{ex}

% nguồn: MaHoaNH_TracNghiem_Lop10
\begin{ex}%[0D2H1-1]%[ma:e4ed2a75]
	Tìm $m$ để  bất phương trình $4x - my < 5$ là bất phương trình bậc nhất hai ẩn?
	\choice 
		{\True  $m \in \mathbb{R}$}
		{$m \ne 0$} 
		 {$m > 0$}
		 {$m < 0$}
	 \loigiai{
		Bất phương trình bậc nhất hai ẩn $x, y$ có  dạng tổng quát là
		\[ax + by \le c,\ ax + by \ge c,\ ax + by < c,\ ax + by  > c \]
		trong đó $a$, $b$, $c$ là những số thực đã cho; $a$ $b$ không đồng thời bằng $0$ và $x$, $y$  là các ẩn.\\
		Do  đó, với mọi số thực $m$ thì bất phương trình $4x - my < 5$ là bất phương trình bậc nhất hai ẩn.
	} 
\end{ex}

% nguồn: MaHoaNH_TracNghiem_Lop10
\begin{ex}%[0D2Y1-1]%[ma:a042df19]
	Điểm nào sau  đây thuộc đồ thị hàm số $y = 2x^2 + x - 3$?
	\choice 
		{\True  $\left(0;-3\right)$}
		{$\left(-2;1\right)$} 
		 {$\left(-1;0\right)$}
		 {$\left(3;-7\right)$}
	\loigiai{Thử tọa độ các  điểm vào hàm số ta được điểm $\left(0;-3\right)$ thuộc đồ thị hàm số.
	} 
\end{ex}

% nguồn: MaHoaNH_TracNghiem_Lop10
\begin{ex}%[0D2B1-1]%[ma:9f12e8d7]
	Lớp $10A$ có $15$ học sinh nam và $20$ học sinh nữ. Hỏi có bao nhiêu cách chọn ra một học sinh để tham gia vào đội thanh niên tình nguyện  của trường, biết rằng tất cả các bạn trong lớp đều có khả năng tham gia.
	 \choice
		 {$\mathrm{C}_{29}^{15}$}
		 {\True $35$}
		{$300$} 
		{$20$} 
	\loigiai{ 
		 Số cách chọn ra $1$ học sinh là $15+20=35$ cách.
	 }
\end{ex}

% nguồn: MaHoaNH_TracNghiem_Lop10
\begin{ex}%[0D2Y1-1]%[ma:a27bfc3c]
	Có $5$ cái bút chì màu khác nhau và $6$ cây bút bi khác nhau. Một bạn học sinh cần chọn $1$ cái bút. Hỏi bạn học sinh  đó có bao nhiêu cách chọn?
	 \choice
		{$30$} 
		 {$6$}
		{\True $11$} 
		{$5$} 
	\loigiai{ 
		Số cách chọn $1$ bút từ $5$ cái bút chì màu khác nhau và $6$ cây bút bi khác nhau  là $5+6=11$.
	} 
\end{ex}

% nguồn: MaHoaNH_TracNghiem_Lop10
\begin{ex}%[0D2Y1-1]%[ma:e2400bf7]
	Lớp 12A có $43$ học sinh, lớp 12B có $30$  học sinh. Hỏi có bao nhiêu cách chọn $1$ học sinh từ lớp 12A và 12B?
	\choice 
		{$43$} 
		 {$30$}
		{\True  $73$}
		{$1290$} 
	 \loigiai{
		Chọn $1$ học sinh từ lớp 12A  và 12B có $43+30=73$ cách.
	} 
\end{ex}

% nguồn: MaHoaNH_TracNghiem_Lop10
\begin{ex}%[0D2B1-1]%[ma:d906f306]
	Bất phương trình nào sau đây là bất phương trình bậc nhất  hai ẩn $x$, $y$?
	 \choice
		 {\True $\sqrt{3}x-y<5^2$}
		{$x^2+3y<4$} 
		 {$\sqrt{2}x-5xy>0$}
		{$x+3y+z-xyz<0$} 
	\loigiai{ 
		Bất phương trình bậc nhất hai ẩn $x$, $y$  là ``$\sqrt{3}x-y<5^2$''.
	 }
\end{ex}

% nguồn: MaHoaNH_TracNghiem_Lop10
\begin{ex}%[0D2Y1-1]%[ma:09232230]
	 Một công việc được hoàn thành bởi một trong hai hành động. Nếu hành động thứ nhất có $m$ cách thực hiện, hành động thứ hai có $n$ cách thực hiện (các cách thực hiện của cả $2$ hành động là khác nhau đôi một) thì số cách hoàn thành công việc đó là
	 \choice
		 {\True $m+n$}
		 {$m-n$}
		{$\dfrac{m}{n}$} 
		 {$m \cdot n$}
	\loigiai{ 
		Số  cách hoàn thành công việc đó là $m+n$.
	} 
\end{ex}

% nguồn: MaHoaNH_TracNghiem_Lop10
\begin{ex}%[0D2Y1-1]%[ma:ab6795fb]
	Cho hàm số $y=f(x)=\heva{&\sqrt{x+4}&&\text{ khi } x>1\\&x^2+1&&\text{ khi } -1\leq x\leq  1\\&2x-1&&\text{ khi } x<-1}$. Giá trị $f(0)$ bằng
	\choice 
		{$-2$} 
		{$2$} 
		{$-1$} 
		{\True $1$} 
	 \loigiai{
		Ta  có $f(0)=0^2+1=1$.
	 }
\end{ex}

% nguồn: MaHoaNH_TracNghiem_Lop10
\begin{ex}%[0D2Y1-1]%[ma:3d02c352]
	Trên giá sách có $5$  cuốn truyện ngắn, $6$ cuốn tiểu thuyết (tất cả đều khác nhau). Có bao nhiêu cách chọn một cuốn để đọc?
	\choice 
		{\True  $11$}
		 {$30$}
		 {$6$}
		 {$5$}
	\loigiai{ 
		Áp dụng quy tắc cộng ta có $5+6=11$ cách chọn một cuốn sách để  đọc.
	 }
\end{ex}

% nguồn: MaHoaNH_TracNghiem_Lop10
\begin{ex}%[0D2Y1-1]%[ma:ebcf2aed]
	Bạn A đi mua một cái áo thu ở một cửa hàng. Giả sử cửa hàng có $10$ cái áo thun size S, $5$ cái áo thun size M, $18$ áo size L.  Giả sử áo size nào bạn A cũng mang được. Hỏi bạn A có bao nhiêu cách lựa chọn?
	\choice 
		 {\True $33$ áo}
		 {$900$ áo}
		 {$15$ áo}
		{$23$ áo} 
	 \loigiai{
		Bạn A có $10+5+18=33$ cách  chọn áo.
	} 
\end{ex}

% nguồn: MaHoaNH_TracNghiem_Lop10
\begin{ex}%[0D2Y1-1]%[ma:8ef6e986]
	Cho hàm số $y=\heva{&2x+1 &\text{khi}\ x\le 2\\ &x^2-3 &\text{khi}\ x>2}$. Trong các điểm sau đây, điểm  nào thuộc đồ thị hàm số?
	 \choice
		{\True  $(0;1)$}
		{$(0;-3)$} 
		{$(3;7)$} 
		 {$(-3;6)$}
	 \loigiai{
		Điểm $(0;1)$ thuộc đồ thị hàm số. 
	} 
\end{ex}

% nguồn: MaHoaNH_TracNghiem_Lop10
\begin{ex}%[0D2Y1-1]%[ma:7f5093a6]
	Giả sử từ tỉnh A đến tỉnh B có thể đi bằng các phương  tiện: ô tô, tàu hỏa, tàu thủy hoặc máy bay.
	Mỗi ngày có $10$ chuyến ô tô, $5$ chuyến tàu hỏa, $3$ chuyến  tàu thủy và $2$ chuyến máy bay. Hỏi có bao nhiêu cách đi từ tỉnh A đến tỉnh B?
	 \choice
		{$20$} 
		{$300$} 
		{\True  $18$}
		 {$15$}
	 \loigiai{
		Đi từ tỉnh A đến tỉnh B có  $10+5+3=18$ cách.
	 }
\end{ex}


%%%%%%%%%%%%%%%% ĐÚNG/SAI %%%%%%%%%%%%%%%%

% nguồn: MaHoaNH_DungSai_Lop10
\begin{ex}%[0D2N1-1]%[ma:c2bb7c6d]
	Cho bất phương trình  $2x-y<3$ (1).
	 \choiceTF
		{\True  Cặp số $(1;1)$ là nghiệm của (1)}
		{ Cặp số  $(2;0)$ là nghiệm của (1)}
		{  Cặp số $(0;-1)$ là nghiệm của (1)}
		{Cặp số $(5;8)$ là nghiệm  của (1)}
	 \loigiai{
		\begin{itemchoice}
			\itemch Cặp số $(1;1)$ là nghiệm  của (1).
			\itemch Cặp số $(2;0)$ không là  nghiệm của (1).
			\itemch Cặp số $(0;-1)$  là nghiệm của (1).
			\itemch Cặp số  $(5;8)$ là nghiệm của (1).
		\end{itemchoice}
	 }
\end{ex}

% nguồn: MaHoaNH_DungSai_Lop10
\begin{ex}%[0D2H1-1]%[ma:8da1d328]
	Xét tính đúng, sai  của các mệnh đề sau
	 \choiceTF
		{\True $\dfrac{-1}{7}x  - \dfrac{y}{3} \leq 8$ là bất phương trình bậc nhất hai ẩn}
		{$\sqrt{2x^2} - 5\sqrt{y}  \geq 8$ là bất phương trình bậc nhất hai ẩn}
		{$2\dfrac{1}{x} -  5\dfrac{1}{y} > 8$ là bất phương trình bậc nhất hai ẩn}
		{\True $\dfrac{2}{5}x - 5^2 y \leq -\sqrt{15}$ là bất  phương trình bậc nhất hai ẩn}
	\loigiai{ 
		\begin{itemchoice}
			\itemch $\dfrac{-1}{7}x - \dfrac{y}{3} \leq 8$ là bất phương trình bậc nhất hai ẩn. 
			\itemch $\sqrt{2x^2} - 5\sqrt{y} \geq 8$ không phải là bất  phương trình bậc nhất hai ẩn.
			\itemch $2\dfrac{1}{x}  - 5\dfrac{1}{y} > 8$ không phải là bất phương trình bậc nhất hai ẩn.
			\itemch $\dfrac{2}{5}x - 25y \leq -\sqrt{15}$  là bất phương trình bậc nhất hai ẩn .
		\end{itemchoice}
	 }
\end{ex}

% nguồn: MaHoaNH_DungSai_Lop10
\begin{ex}%[0D2H1-1]%[ma:dd967d0b]
	Một đội sản xuất cần $3$  giờ để làm xong sản phẩm loại I và $2$ giờ để làm xong sản phẩm loại II. Biết thời gian tối đa cho việc sản xuất hai loại sản phẩm trên là $18$ giờ. Gọi $x$, $y$ lần lượt là số sản phẩm loại I, II mà đội làm được trong thời gian cho phép.
	 \choiceTF
		{  Thời gian làm xong $x$ sản phẩm loại I là $2x$, thời gian làm xong $y$ sản phẩm loại II là $3y$}
		{ Thời gian sản  xuất hai sản phẩm loại I và II thỏa mãn bất phương trình là $3x + 2y < 18$, $x$, $y \in \mathbb{N}$}
		{\True $(3;4)$ là một nghiệm của bất phương trình bậc nhất hai ẩn $3x + 2y \leq  18$ với mọi $x$, $y \in \mathbb{N}$}
		{\True $(4;3)$ là một nghiệm của bất phương  trình bậc nhất hai ẩn $3x + 2y \leq 18$ với mọi $x$, $y \in \mathbb{N}$}
	\loigiai{ 
		\begin{itemchoice}
			\itemch Tổng thời gian làm xong sản phẩm loại I  là $3x$, loại II là $2y$.
			\itemch Ta có bất phương trình $3x  + 2y \leq 18 \quad (*)$ với điều kiện $x$, $y \in \mathbb{N}$.
			\itemch Thay $(3;4)$ vào $(*)$ ta có $3 \cdot 3 + 2 \cdot 4 = 9  + 8 = 17 \leq 18$ (đúng), nên $(3;4)$ là nghiệm.
			\itemch Thay $(4;3)$ vào $(*)$ ta có $3 \cdot 4 + 2 \cdot 3 =  12 + 6 = 18 \leq 18$ (đúng), nên $(4;3)$ là nghiệm của $(*)$.
		\end{itemchoice}
	 }
\end{ex}

% nguồn: MaHoaNH_DungSai_Lop10
\begin{ex}%[0D2H1-1]%[ma:250f62c2]
	An thích ăn hai loại trái cây  là cam và xoài, mỗi tuần mẹ cho An $200\,000$ đồng để mua trái cây. Biết rằng giá cam là $15\,000$ đồng/$1$ kg, giá xoài là $30\,000$ đồng/$1$ kg. Gọi $x$, $y$ lần lượt là số kilogam cam và xoài mà An có thể mua về sử dụng trong một tuần.
	\choiceTF 
		{Trong tuần, số tiền An mua cam là $15\,000x$ đồng, số tiền An có thể mua  xoài là $30\,000y$ đồng $(x$, $y > 0)$}
		{ Bất phương trình $3x + 6y \geq 40$ thể hiện số tiền An có thể mua cam và  xoài trong một tuần}
		{\True Cặp số $(5; 4)$ thỏa mãn bất phương trình $3x + 6y \geq 40$} 
		{ An có thể mua $4$ kg cam, $5$ kg xoài trong tuần} 
	\loigiai{ 
		\begin{itemchoice}
			\itemch Trong tuần, số tiền An có thể mua cam là $15\,000x$, xoài là $30\,000y$ với $(x,  y \geq 0)$.
			\itemch Ta có bất phương trình $15\,000x + 30\,000y \leq 200\,000  \Leftrightarrow 3x + 6y \leq 40. \quad (*)$
			\itemch Thay $x = 5$, $y = 4$ vào (*) ta có $3\cdot5 + 6\cdot4 =  15 + 24 = 39 \leq 40$ (đúng) nên $(5;4)$ là nghiệm của $3x + 6y \leq 40$.
			\itemch Thay $x = 4$, $y = 5$ vào (*) ta có $3\cdot4 + 6\cdot5 = 42 > 40$ nên An không mua được $4$ kg cam, $5$ kg xoài trong  tuần.
		\end{itemchoice}
	} 
\end{ex}

% nguồn: MaHoaNH_DungSai_Lop10
\begin{ex}%[0D2H1-1]%[ma:1b72b08a]
	Một trò chơi chọn ô chữ đơn giản mà kết quả gồm một trong hai khả năng: Nếu người chơi chọn được chữ $A$ thì người ấy được cộng $3$ điểm, nếu người chơi chọn được chữ $B$ thì người ấy bị  trừ $1$ điểm. Người chơi chỉ chiến thắng khi đạt được số điểm tối thiểu là $20$. Gọi $x$, $y$ theo thứ tự là số lần người chơi chọn được chữ $A$ và chữ $B$. Khi đó:
	\choiceTF 
		 {\True Tổng số điểm người chơi đạt được khi chọn chữ $A$ là $3x$, tổng số điểm người chơi bị trừ khi chọn chữ $B$ là $y$}
		{Bất phương trình bậc nhất hai ẩn trong tình huống người chơi chiến thắng  là $3x - y \geq 18$}
		{\True Người chơi chọn được chữ $A$ $7$ lần và chọn được chữ $B$ $1$ lần thì người đó vừa  đủ điểm dành chiến thắng trò chơi}
		{\True Người  chơi chọn được chữ $A$ $8$ lần và chọn được chữ $B$ $3$ lần thì người đó vừa đủ điểm dành chiến thắng trò chơi}
	\loigiai{ 
		\begin{itemchoice}
			\itemch Tổng điểm đạt được là $3x$, bị trừ là $y$ nên tổng điểm còn lại  là $3x - y$.
			\itemch Điều kiện là $3x  - y \geq 20$ mới đủ điểm chiến thắng.
			\itemch Với $x =  7$, $y = 1$ thì $3x - y = 21 - 1 = 20 \geq 20$ (đúng).
			\itemch Với $x  = 8$, $y = 4$ thì $3x - y = 24 - 4 = 20 \geq 20$ (đúng).
		\end{itemchoice}
	 }
\end{ex}


%%%%%%%%%%%%%%%% TRẢ LỜI NGẮN %%%%%%%%%%%%%%%%

% nguồn: MaHoaNH_TraLoiNgan_Lop10
\begin{bt}%[0D2H1-1]%[ma:fc55348e]
	Bạn Việt mang $100\,000$ đồng ra chợ mua hoa cúc và hoa hồng. Một bông hoa cúc có giá $3\,000$ đồng, một bông hoa hồng có giá $6\,000$ đồng. Gọi $x$ và $y$ lần lượt là số bông hoa cúc và số bông hoa hồng  bạn Việt mua. Bất phương trình bậc nhất hai ẩn $x$, $y$ để biểu diễn số tiền Việt mua hoa cúc và hoa hồng có dạng $ax + 6y \leq b$ với $a, b \in \mathbb{N}$. Tính giá trị biểu thức $T = a + b$.
	\shortans{103} 
	 \loigiai{
		Một bông hoa cúc có giá $3\,000$ đồng, mua $x$ bông  hết $3\,000x$ đồng.\\
		Một bông hoa hồng có giá $6\,000$ đồng, mua $y$ bông hết $6\,000y$  đồng.\\
		Bạn Việt mang $100\,000$ đồng ra chợ mua hoa cúc và hoa hồng nên ta có  điều kiện:
		\[ 
		3\,000x + 6\,000y \leq 100\,000 \Leftrightarrow  3x + 6y \leq 100.
		 \]
		Vậy $a = 3$, $b = 100 \Rightarrow T =  a + b = 103$.
	 }
\end{bt}

% nguồn: MaHoaNH_TraLoiNgan_Lop10
\begin{bt}%[0D2H1-1]%[ma:caacb4d2]
	Bạn Hương mang $600\,000 $ đồng đi siêu thị mua thực phẩm (gồm thịt và rau) cho gia đình gồm $4$ người dùng trong $4$ ngày. Biết rằng mỗi kg thịt có giá $120\,000$ đồng, mỗi kg rau có giá $30\, 000$ đồng và siêu  thị chỉ bán hàng theo kg chứ không bán lẻ. Nếu gọi $x$ là số kg thịt và $y$ là số kg rau mà Hương mua thì điều kiện của $x$, $y$ thỏa mãn bất phương trình $ax + by \le 20$ với $a$, $b \in \mathbb{N}$.
	Tính giá trị biểu thức $T = 2a +  3b$.
	 \shortans{11}
	 \loigiai{
		Vì cửa hàng không bán lẻ  nên $x, y \in \mathbb{N}^*$.\\
		 Số tiền mà Hương mua thịt là $120\,000x$ đồng và số tiền mà Hương mua rau là $30\,000y$ đồng.\\
		Vì Hương chỉ có 600 000 đồng nên ta có bất  phương trình
		\[ 
		120\,000x + 30\,000y \le  600\,000 \Leftrightarrow 4x + y \le 20.
		\] 
		Vậy $a = 4$, $b =  1 \Rightarrow T = 2a + 3b = 2\cdot 4 + 3\cdot1 = 11$.
	 }
\end{bt}

% nguồn: MaHoaNH_TraLoiNgan_Lop10
\begin{bt}%[0D2H1-1]%[ma:5705970c]
	Một cửa hàng bán hai loại đồ uống có tên là \lq\lq Giọt lệ thiên thần\rq\rq \ và \lq\lq Giọt lệ ác quỷ\rq\rq. Bốn ly \lq\lq Giọt lệ thiên thần\rq\rq \ có giá $600\,000$ đồng, ba ly \lq\lq Giọt lệ ác quỷ\rq\rq \ có giá $540\,000$ đồng. Hàng tháng, cửa hàng này phải chi trả $6\,000\,000$ đồng tiền thuê nhân viên, $8\,000\,000$ đồng tiền thuê mặt bằng, $3\,000\,000$ đồng tiền nguyên liệu. (Ngoài ra cửa hàng không tốn thêm bất kỳ chi phí gì và thu nhập của cửa hàng chỉ đến từ việc bán hai loại đồ uống trên). Gọi $x$, $y$ lần lượt là số ly \lq\lq Giọt lệ thiên thần\rq\rq \ và \lq\lq Giọt lệ ác quỷ\rq\rq \  mà cửa hàng bán được trong một tháng. Điều kiện của $x$, $y$ để doanh thu của cửa hàng trong một tháng có lãi thoả mãn bất phương trình $ax + by \leq 1\,700$ với $a$, $b \in \mathbb{N}$. Tính giá trị biểu thức $T = 2a + b$.
	 \shortans{48}
	\loigiai{ 
		\begin{itemize}
			\item $4$ ly \lq\lq Giọt lệ thiên thần\rq\rq\  có giá $600\,000$ đồng nên 1 ly có giá $150\,000$ đồng.
			\item $3$ ly \lq\lq Giọt lệ ác quỷ\rq\rq\ có giá $540\,000$ đồng nên 1 ly có  giá $180\,000$ đồng.
			\item Tổng chi phí trong tháng là $17\,000\,000$ đồng. 
			\item Để cửa hàng có lãi thì  doanh thu cần lớn hơn $17\,000\,000$ đồng, tức là
			\[ 
			150\,000x + 180\,000y > 17\,000\,000 \Leftrightarrow 15x  + 18y > 1\,700.
			 \]
		\end{itemize}
		Vậy $a = 15$, $b = 18 \Rightarrow  T = 2a + b = 2 \cdot 15 + 18 = 48$.
	 }
\end{bt}

% nguồn: MaHoaNH_TraLoiNgan_Lop10
\begin{bt}%[0D2H1-1]%[ma:7defc288]
	Cho bất phương trình $x + 3y - 12 \geq 0$. Có bao nhiêu giá trị nguyên của tham số $m$ để cặp số $(m^2; m^2 + 2m - 2)$ không phải là nghiệm của bất phương trình đã cho? 
	 \shortans{4}
	\loigiai{ 
		Do cặp $(m^2; m^2 + 2m - 2)$ không là  nghiệm của bất phương trình $x + 3y - 12 \geq 0$ nên ta có
		\begin{align*}
			& m^2 + 3(m^2 + 2m - 2)  - 12 < 0 \\
			& \Leftrightarrow m^2 + 3m^2 + 6m  - 6 - 12 < 0 \\
			& \Leftrightarrow 4m^2  + 6m - 18 < 0 \Leftrightarrow (2m - 3)(m + 3) < 0.
		\end{align*}
		Xét  dấu biểu thức:
		 \[
		(2m  - 3)(m + 3) < 0 \Leftrightarrow
		\hoac{& \heva{& 
		2m  - 3 > 0 \\&
		m + 3 < 0 }\\&\heva{&2m-3>0\\&m+3<0}} 
		\Leftrightarrow \hoac{&\heva{&m > \dfrac{3}{2}\\&m < -3}\\&\heva{&m<\dfrac{3}{2}\\&m>-3}}\Leftrightarrow  -3<m<\dfrac{3}{2}.
		 \]
		 Vì
		 $
		-3 < m <  \dfrac{3}{2} \quad \text{và } m \in \mathbb{Z} \Rightarrow m \in \{-2; -1; 0; 1\}.
		 $
		Vậy có tất cả $4$ giá trị nguyên của tham số  $m$ thỏa mãn.
	} 
\end{bt}

% nguồn: MaHoaNH_TraLoiNgan_Lop10
\begin{bt}%[0D2V1-1]%[ma:8a5ec331]
	Cho bất phương trình $2x + 3y - 10 \leq 0$. Hỏi có bao nhiêu cặp số nguyên $(m_0; n_0)$ thỏa mãn $(m_0^2; n_0^2)$ là nghiệm của bất phương  trình đã cho?
	\shortans{11} 
	\loigiai{ 
		Vì $(m_0^2;  n_0^2)$ là nghiệm của bất phương trình $2x + 3y - 10 \leq 0$ nên ta có
		\[ 
		2m_0^2 + 3n_0^2  \leq 10 \Rightarrow
		 \heva{&
		 m_0^2 \leq 5 \\&
		n_0^2 \leq  \dfrac{10}{3} }
		\Leftrightarrow 
		\heva{&-\sqrt{5} \leq m_0 \leq \sqrt{5}\\&-\sqrt{\dfrac{10}{3}}  \leq n_0 \leq \sqrt{\dfrac{10}{3}}.}
		\] 
		Do $m_0, n_0  \in \mathbb{Z}$ nên:
		\[ 
		m_0 \in \{-2; -1; 0; 1; 2\}, \quad n_0 \in \{-1;  0; 1\}.
		\] 
		Loại các cặp không thỏa: $(-2; -1)$, $(2; 1)$, $(2; -1)$, $(-2; 1)$, còn lại $11$ cặp  thỏa mãn.
		Vậy có $11$ cặp số $(m_0; n_0)$ sao cho $(m_0^2; n_0^2)$ là nghiệm của bất phương trình đã  cho.
	 }
\end{bt}


%%%%%%%%%%%%%%%% TỰ LUẬN %%%%%%%%%%%%%%%%

% nguồn: MaHoaNH_TuLuan_Lop10
\begin{bt}%[0D2N1-1]%[ma:9daca6c3]
	 Tìm các bất phương trình bậc nhất hai ẩn trong các bất phương trình sau
	\begin{multicols}{2}
		\begin{enumerate}
			\item $x-5y+2  \le 0$.
			 \item $9x^2+8y-7 \ge 0$.
			 \item $3x-2y>0$.
			\item $4y+11x  \le 0$.
		\end{enumerate}
	\end{multicols}
	\loigiai{Dựa trên định nghĩa ta thấy bất phương trình bậc nhất hai ẩn là các bất phương  trình a), c), d).\\
		Bất phương trình b) không phải là bất  phương trình bậc nhất hai ẩn vì có chứa $x^2$.
	 }
\end{bt}

% nguồn: MaHoaNH_TuLuan_Lop10
\begin{bt}%[0D2N1-1]%[ma:5e70b011]
	Tìm bất  phương trình bậc nhất hai ẩn trong các bất phương trình sau
	\begin{multicols}{3}
		\begin{enumerate}
			\item $3x-2(y-x+1)>0$. 
			 \item $2^2x+y \le 0$.
			\item $2x^2-y  \ge 1$.
		\end{enumerate}
	\end{multicols}
	 \loigiai{
		Ta có $3x-2(y-x+1) >0 \Leftrightarrow 3x-2y+2x-2>0  \Leftrightarrow 5x-2y-2>0$.\\
		Và $2^2x+y \le  0 \Leftrightarrow 4x+y \le 0$.\\
		Vậy bất phương trình bậc nhất hai ẩn là  hai bất phương trình ở câu a) và b).\\
	Bất phương trình $2x^2-y \ge 1$ không phải là bất phương trình bậc  nhất hai ẩn vì chứa $x^2$.}
\end{bt}

% nguồn: MaHoaNH_TuLuan_Lop10
\begin{bt}%[0D2H1-1]%[ma:fd7ea1ec]
	Tìm $m$ để bất phương  trình sau là bất phương trình bậc nhất hai ẩn
	 $$\left(m^2+m\right)x^2-mx+m^2y-1 \le 0.$$
	\loigiai{Bất phương trình đã cho là bất phương trình bậc nhất hai ẩn khi và chỉ khi 
		$$\heva{&m^2+m=0\\&(-m)^2+(m^2)^2 \neq 0} \Leftrightarrow \heva{&\hoac{&m=0\\&m=-1}\\&m^2 \left(1+m^2\right) \neq  0} \Leftrightarrow \heva{&\hoac{&m=0\\&m=-1}\\&m \neq 0} \Leftrightarrow m=-1.$$
	 Vậy $m=-1$.}
\end{bt}

% nguồn: MaHoaNH_TuLuan_Lop10
\begin{bt}%[0D2T1-1]%[ma:7766f340]
	Một công ty viễn thông A cung cấp dịch vụ truyền hình cáp với mức phí ban đầu là 300000 đồng và mỗi tháng phải đóng 150000 đồng. Công ty viễn thông B cũng cung cấp dịch vụ truyền hình cáp nhưng không tính phí ban  đầu và mỗi tháng khách hàng sẽ phải đóng 200000 đồng.
	\begin{tasks}(1)
		\task Gọi  $T$ (đồng) là số tiền khách hàng phải trả cho mỗi công ty viễn thông trong $t$ (tháng) sử dụng dịch vụ truyền hình cáp. Khi đó hãy lập hàm số $T$ theo $t$ đối với mỗi công ty.
		\task Tính số tiền khách hàng phải trả sau khi sử dụng dịch vụ truyền hình cáp trong 5 tháng đối  với mỗi công ty.
		\task Khách  hàng cần sử dụng dịch vụ truyền hình cáp trên mấy tháng thì đăng kí bên công ty viễn thông A sẽ tiết kiệm chi phí hơn?
	\end{tasks}
	 \loigiai{
		\begin{enumerate}[a)]
			\item Hàm số $T$ theo $t$ đối với công  ty A là $T=150000t+300000$.\\
			 Hàm số $T$ theo $t$ đối với công ty B là $T=200000t$.
			\item Thay $t=5$ lần lượt vào  hai công thức trên, ta được
			\begin{itemize}
				\item [$\bullet$] Số tiền phải trả trong 5 tháng khi sử dụng dịch  vụ truyền hình cáp của công ty A là $150000.5+300000=1050000$ đồng.
				\item [$\bullet$] Số tiền phải trả trong 5 tháng khi sử dụng  dịch vụ truyền hình cáp của công ty B là $200000.5=1000000$ đồng.
			\end{itemize}
			 \item Để dịch vụ truyền hình cáp của công ty A lợi hơn dịch vụ truyền hình cáp của công ty B thì:
			$$ 
			150000t+300000<200000t \Leftrightarrow 300000<50000t  \Leftrightarrow t>6
			$$ 
			Vậy nếu sử dụng từ 7 tháng trở lên thì sử dụng dịch vụ truyền hình cáp bên công ty A sẽ có lợi  hơn.
	\end{enumerate}}
\end{bt}

% nguồn: MaHoaNH_TuLuan_Lop10
\begin{bt}%[0D2B1-1]%[ma:2058088d]
	Sau khi đun nóng băng phiến lên đến gần $90^{\circ}C$, người ta để nguội, quan sát, ghi nhận nhiệt độ và trạng thái của băng phiến sau mẫu phút như bảng  sau
	\begin{center}
		\textit{Nhiệt độ và trạng thái của băng phiến khi để  nguội}
		\begin{tabular}{|l|c|c|c|c|c|c|c|c|c|c|c|}
			\hline 
			\textbf{Thời gian nguội (phút)} & $0$ & $1$ & $2$ & $3$ & $4$ & $5$ & $6$ & $7$ & $8$ & $9$ & $10$ \\ 
			 \hline
			\textbf{Nhiệt độ ($^{\circ}C$)} & $86$ & $84$ & $82$ &  $81$ & $80$ & $80$ & $80$ &$80$ & $79$ & $77$ & $75$ \\
			\hline 
			\textbf{Trạng thái} & \multicolumn{4}{|c|}{lỏng} &\multicolumn{4}{|c|}{lỏng và rắng} & \multicolumn{3}{|c|}{rắn}\\ 
			 \hline
		\end{tabular}
	\end{center}
	\begin{tasks}(1)
		\task Tại sao từ  bảng trên, có thể nói nhiệt độ của băng phiến là một hàm số theo thời gian (nung nóng)? Tìm tập xác định và tập giá trị của hàm số trên.
		\task Sau khi  để nguội $3$ phút, nhiệt độ băng phiến là bao nhiêu?
		\task Băng phiến chuyển hoàn toàn sang trạng thái rắn sau  bao nhiêu phút?
	\end{tasks}
	\loigiai{ 
		\begin{enumerate}[a)]
			\item Bảng giá trị cho thấy  nhiệt độ (kí hiệu là $y$) là một hàm số theo thời gian (kí hiệu là $x$) vì khi cho $x$ một giá trị bất kì, ta luôn tìm được duy nhất một giá trị của $y$. Do vậy bảng này xác định một hàm số biểu thị nhiệt độ của băng phiến theo thời gian.\\
			Từ bảng giá trị của hàm số, ta có tập xác định $\mathscr{D}=\{0; 1; 2;  3; 4; 5; 6; 7; 8; 9; 10\}$ và tập giá trị $T=\{75; 77; 79; 80; 81; 82; 84; 86\}$.
			\item Sau khi để nguội  $3$ phút, nhiệt độ băng phiến là $81^{\circ}C$.
			\item Băng phiến chuyển hoàn toàn sang trạng thái rắn sau $8$ phút (lúc đó nhiệt độ băng phiến  là $79^{\circ}$).
		\end{enumerate}
	 }
\end{bt}

% nguồn: MaHoaNH_TuLuan_Lop10
\begin{bt}%[0D2T1-1]%[ma:23e251a5]
	Theo quyết định số 2019/QĐ-BĐVN  ngày 01/11/2018 của Tổng công ty Bưu điện Việt Nam, giá cước dịch vụ Bưu chính phổ cập đối với dịch vụ thư cơ bản và bưu thiếp trong nước có khối lượng đến $250$g như trong bảng sau
	\immini{\begin{enumerate}[a)]
	\item Số tiền dịch vụ thư cơ bản phải trả $y$ (đồng) có là hàm số  của khối lượng thư cơ bản $x$ (g) hay không? Nếu đúng, hãy xác định những công thức tính $y$.
	\item Tính số tiền phải trả khi bạn Dương gửi thư có khối lượng $150$g, $200$g. 
\end{enumerate} }
 {\vspace{0.6 cm}
\begin{tabular}{|c|c|}
	 \hline
	\textcolor{blue}{Khối lượng  đến $250$ g} & \textcolor{blue}{Mức cước (đồng)}\\
	 \hline
	 Đến $20$ g & $4000$\\
	\hline 
	 Trên $20$ g đến $100$ g & $6000$ \\
	 \hline
	Trên $100$ g  đến $250$ g & $8000$\\
	 \hline
\end{tabular}}
\loigiai{ 
	\begin{enumerate}[a)]
		\item Số tiền dịch vụ thư cơ bản phải  trả $y$ là hàm số của $x$ vì mỗi giá trị của $x$ (chính là khối lượng của thư) có đúng một giá trị của $y$ (mức cước hay số tiền phải trả) tương ứng.\\
		Quan  sát bảng ta thấy:
		\begin{itemize}
			\item[$\bullet$] Nếu  khối lượng thư đến $20$ g hay $0<x \le 20$ thì mức cước phải trả là $4000$ đồng hay $y= 4000$.
			\item[$\bullet$] Nếu khối lượng thư trên $20$ g hay $100$ g hay $20 <  x \le 100$ g thì mức cước là $6000$ đồng hay $y=6000$.
			\item[$\bullet$] Nếu khối lượng thư trên $100$ g đến $250$ g hay $100 < x \le 250$ thì  mức cước là $8000$ đồng hay $y=8000$.
		\end{itemize}
		 Vậy ta có công thức xác định $y$ như sau:
		$y=\heva{&4000&\text{nếu}& 0<x\le 20\\&6000&\text{nếu}& 20 <x\le 100\\&8000&\text{nếu}& 100<x\le 250.}$ 
		\item Vì $100 < 150 <250$ và $100 <200  <250$ nên bức thư có khối lượng $150$ g thì cần trả cước là $8000$ đồng và bức thư có khối lượng $200$ g cũng cần trả cước là $8000 $ đồng. \\
		Vậy  tổng số tiền phải trả khi bạn Dương gửi thư có khối lượng $150$ g và $200$ g là \\ \centerline{$8000 + 8000 = 16 000$ (đồng).}
	\end{enumerate}
 }
\end{bt}
